\title{The Era of 1-bit LLMs: All Large Language Models are in 1.58 Bits}

\begin{abstract}
Recent advancements, including BitNet~\cite{bitnet}, are ushering in a new period for 1-bit Large Language Models (LLMs). In this study, we propose a 1-bit LLM variant called \textbf{\text{BitNet b1.58}}, where every parameter (or weight) of the LLM is ternarized to \{-1, 0, 1\}. This model achieves comparable perplexity and downstream task performance to full-precision (i.e., FP16 or BF16) Transformer LLMs of the same size and trained on equivalent token counts, while offering substantial improvements in latency, memory usage, throughput, and energy efficiency. More importantly, the 1.58-bit LLM establishes a novel scaling law and training methodology for developing future LLMs that are both efficient and high-performing. Additionally, it introduces a new computational paradigm and lays the foundation for specialized hardware designs tailored to 1-bit LLMs.
\end{abstract}

\section{The Era of 1-bit LLMs}

Over the past few years, the AI community has witnessed rapid advancements in the scale and capabilities of Large Language Models (LLMs). These models have excelled across numerous natural language processing tasks, yet their growing sizes present deployment difficulties and raise environmental and economic concerns due to significant energy demands.

One strategy to mitigate these issues is post-training quantization, which produces low-bit models optimized for inference~\cite{smoothquant, gptq, quip, quip_sharp}. This approach lowers the precision of weights and activations, dramatically cutting down memory and computational needs for LLMs. The progression has typically involved decreasing bit widths from 16 bits down to variants with as low as 4 bits~\cite{gptq, awq}. Nonetheless, despite its widespread industry use, post-training quantization remains suboptimal.

Recent research on 1-bit model architectures, such as BitNet~\cite{bitnet}, offers a promising approach to decrease the cost of LLMs while preserving their performance. Standard LLMs operate with 16-bit floating-point values (i.e., FP16 or BF16), and the primary operation in any LLM is matrix multiplication. Consequently, the main computational expense arises from floating-point addition and multiplication operations. In contrast, BitNet’s matrix multiplication uses only integer addition, which drastically reduces the energy consumption for LLMs. Since power is often the limiting factor for compute performance in many chips, these energy savings can also translate to faster processing.

Beyond computation, transferring model parameters from DRAM to the memory of an on-chip accelerator (such as SRAM) can be costly during inference. Attempts to increase SRAM size to boost throughput have been made, but this comes with significantly higher costs compared to DRAM. Relative to full-precision models, 1-bit LLMs have a considerably smaller memory footprint in terms of both capacity and bandwidth. This reduction can greatly lower the cost and duration of loading weights from DRAM, resulting in faster and more efficient inference.

In this paper, we propose a notable 1-bit LLM variant termed \textbf{\text{BitNet b1.58}}, where every parameter is ternary, taking values from \{-1, 0, 1\}. By introducing the value 0 to the original 1-bit BitNet, we achieve 1.58 bits in the binary system. \text{BitNet b1.58} preserves all advantages of the original 1-bit BitNet, including its novel computation approach that almost eliminates multiplication operations for matrix multiplication and enables high optimization. It also maintains the same energy consumption as the original 1-bit BitNet, while being substantially more efficient in memory usage, throughput, and latency compared to FP16 LLM baselines. Moreover, \text{BitNet b1.58} provides two additional benefits. First, its modeling capability is enhanced due to explicit feature filtering enabled by the inclusion of 0 in the model weights, which can significantly boost the performance of 1-bit LLMs. Second, our experiments demonstrate that \text{BitNet b1.58} matches full precision (i.e., FP16) baselines in perplexity and end-task performance starting from a 3B parameter size, when using the same configuration (e.g., model size, training tokens, etc.).

\section{BitNet b1.58}

\text{BitNet b1.58} is derived from the BitNet architecture, which is a Transformer that substitutes \emph{nn.Linear} layers with \emph{BitLinear}. It is trained from the beginning, utilizing 1.58-bit weights and 8-bit activations. Compared to the original BitNet, it incorporates several changes, which we outline below.

\paragraph{Quantization Function.} To restrict the weights to the values -1, 0, or +1, we use an \emph{absmean} quantization function. This function first normalizes the weight matrix by its average absolute value, then rounds each element to the nearest integer within the set \{-1, 0, +1\}:

\begin{equation}
    \widetilde{W} = \mathrm{RoundClip}\left(\frac{W}{\gamma + \epsilon}, -1, 1\right),
\end{equation}

\begin{equation}
    \mathrm{RoundClip}(x, a, b) = \max\left(a, \min\left(b, \mathrm{round}(x)\right)\right),
\end{equation}

\begin{equation}
    \gamma = \frac{1}{nm}\sum_{ij} |W_{ij}|.
\end{equation}

For activations, the quantization function follows the same approach as in BitNet, except that we do not scale activations before applying the nonlinear functions into the range $[0, Q_b]$. Instead, activations are scaled per token to the interval $[-Q_b, Q_b]$ to eliminate zero-point quantization. This approach simplifies implementation and system-level optimization, and our experiments show it has minimal impact on performance.

\paragraph{LLaMA-alike Components.}

The LLaMA~\cite{llama, llama2} architecture has become the standard backbone for open-source large language models. To align with the open-source ecosystem, \text{BitNet b1.58} adopts LLaMA-like components. Specifically, it employs RMSNorm~\cite{rmsnorm}, SwiGLU~\cite{swiglu}, rotary embeddings~\cite{rope}, and removes all biases. This design enables \text{BitNet b1.58} to be easily integrated into widely-used open-source frameworks (e.g., Huggingface, vLLM~\cite{vllm}, and llama.cpp\footnote{\url{https://github.com/ggerganov/llama.cpp}}) with minimal effort.

\section{Results}

We conducted a comparison between \text{BitNet b1.58} and our replicated FP16 \text{LLaMA LLM} across various model sizes. To maintain fairness, both models were pre-trained on the RedPajama dataset~\cite{redpajama} with 100 billion tokens. We assessed their zero-shot performance on a diverse set of language tasks, including ARC-Easy~\cite{arc}, ARC-Challenge~\cite{arc}, Hellaswag~\cite{hellaswag}, Winogrande~\cite{winoGrande}, PIQA~\cite{piqa}, OpenbookQA~\cite{openbookqa}, and BoolQ~\cite{boolq}. Additionally, we reported validation perplexity on the WikiText2~\cite{wikitext2} and C4~\cite{c4} datasets.

We also compared the runtime GPU memory usage and latency between \text{LLaMA LLM} and \text{BitNet b1.58}. These metrics were obtained using the FasterTransformer\footnote{\url{https://github.com/NVIDIA/FasterTransformer}} codebase, which is highly optimized for LLM inference latency on GPUs. The 2-bit kernel from Ladder~\cite{ladder} was incorporated for \text{BitNet b1.58}. Since inference cost is largely dominated by token generation time, we reported the time taken per output token.

Table~\ref{tab:ppl} presents a summary of perplexity and computational cost for both \text{BitNet b1.58} and \text{LLaMA LLM}. The data indicate that at the 3B model scale, \text{BitNet b1.58} begins to rival the full precision \text{LLaMA LLM} in terms of perplexity, while achieving a 2.71-fold speedup and requiring 3.55 times less GPU memory. Specifically, the 3.9B version of \text{BitNet b1.58} runs 2.4 times faster, utilizes 3.32 times less memory, and substantially outperforms the \text{LLaMA LLM} 3B model.

Detailed zero-shot accuracy results on downstream tasks are listed in Table~\ref{tab:zero-shot}. We followed the evaluation pipeline provided by \emph{lm-evaluation-harness}\footnote{\url{https://github.com/EleutherAI/lm-evaluation-harness}}. The findings reveal that the performance gap between \text{BitNet b1.58} and \text{LLaMA LLM} shrinks with increasing model size. Crucially, starting at 3B, \text{BitNet b1.58} can match the performance of the full precision baseline. In line with the perplexity results, the zero-shot task outcomes demonstrate that \text{BitNet b1.58} at 3.9B surpasses \text{LLaMA LLM} 3B while incurring lower memory and latency costs. This establishes \text{BitNet b1.58} as a Pareto improvement relative to leading LLM models.

\paragraph{Memory and Latency}

We additionally increased the model sizes to 7B, 13B, and 70B and assessed the associated costs. Figure~\ref{fig:latency-memory-energy} displays the latency and memory trends, indicating that the speed-up improves as the model size increases. Specifically, \text{BitNet b1.58} 70B runs 4.1 times faster than the \text{LLaMA LLM} baseline. This improvement arises because the time required for \emph{nn.Linear} operations grows with the size of the model. Memory usage exhibits a comparable pattern, since the embedding remains in full precision, and its share of memory decreases for larger models. Both latency and memory measurements were performed using a 2-bit kernel, suggesting further optimization potential to lower costs even more.

\paragraph{Energy}

We also estimated the energy consumption of arithmetic operations for both \text{BitNet b1.58} and \text{LLaMA LLM}. Our focus was primarily on matrix multiplication calculations, as they dominate the computational expense of LLMs. Figure~\ref{fig:energy} presents the breakdown of energy costs. Most of \text{BitNet b1.58}'s energy use comes from INT8 addition operations, whereas \text{LLaMA LLM} consumes energy through both FP16 addition and FP16 multiplication. Based on the energy model described in~\cite{energycost, pokebnn}, \text{BitNet b1.58} reduces arithmetic operation energy consumption for matrix multiplication by 71.4 times on 7nm chips. We also reported the total energy cost for models processing 512 tokens. Our findings indicate that as model size increases, \text{BitNet b1.58} becomes progressively more energy efficient compared to the FP16 \text{LLaMA LLM} baseline. This efficiency gain is attributed to the growing proportion of \emph{nn.Linear} operations with larger models, while the relative cost of other components decreases.

\paragraph{Throughput}

We evaluate the throughput of \text{BitNet b1.58} and \text{LLaMA LLM} with 70B parameters using two 80GB A100 GPUs, employing pipeline parallelism~\cite{gpipe} to enable running the \text{LLaMA LLM} 70B on these devices. The batch size was progressively increased until the GPU memory limit was reached, with a fixed sequence length of 512. As shown in Table~\ref{tab:throughput}, \text{BitNet b1.58} 70B can accommodate a batch size up to 11 times larger than that of \text{LLaMA LLM}, resulting in a throughput that is 8.9 times greater.

\textbf{\text{BitNet b1.58} establishes a new scaling relationship between model performance and inference cost}. For reference, based on the outcomes presented in Figures~\ref{fig:latency-memory-energy} and \ref{fig:energy}, the following equivalences hold between various model sizes in 1.58-bit and 16-bit precision:

\begin{itemize}
    \item A 13B \text{BitNet b1.58} model is more efficient in latency, memory consumption, and energy usage than a 3B FP16 LLM.
    \item A 30B \text{BitNet b1.58} model outperforms a 7B FP16 LLM in terms of latency, memory utilization, and energy efficiency.
    \item A 70B \text{BitNet b1.58} model is more efficient regarding latency, memory usage, and energy consumption than a 13B FP16 LLM.
\end{itemize}

\paragraph{Training with 2T Tokens}

The amount of training tokens plays a vital role in the effectiveness of LLMs. To assess the scalability of \text{BitNet b1.58} concerning token volume, we trained a \text{BitNet b1.58} model using 2 trillion tokens, following the data preparation protocol of StableLM-3B~\cite{StableLM-3B-4E1T}, which is a leading open-source 3B model. Both models were evaluated on a benchmark composed of Winogrande~\cite{winoGrande}, PIQA~\cite{piqa}, SciQ~\cite{sciq}, LAMBADA~\cite{lambada}, and ARC-easy~\cite{arc}. Zero-shot accuracy results are presented in Table~\ref{tab:2t}. For tasks assessed by accuracy and normalized accuracy, the average of both metrics was taken. The results for StableLM 3b at 2T tokens were sourced from its technical report. Our results demonstrate that \text{BitNet b1.58} achieves superior performance across all tasks, indicating that 1.58-bit LLMs also possess strong generalization abilities.

\section{Discussion and Future Work}

\textbf{1-bit Mixture-of-Experts (MoE) LLMs}

Mixture-of-Experts (MoE) models have shown to be an efficient solution for large language models (LLMs) by substantially lowering computation FLOPs. However, their deployment and usability are hindered by high memory demands and the communication overhead between chips. These issues can be mitigated by using 1.58-bit LLMs. Primarily, the diminished memory requirement decreases the number of devices needed for MoE model deployment. Additionally, it greatly cuts down the cost of transferring activations over networks. Ultimately, if the entire model fits onto a single chip, this overhead would be eliminated entirely.

\textbf{Native Support of Long Sequence in LLMs}

With the rise of LLMs, the capacity to process long sequences has become essential. A key obstacle in long sequence inference is the substantial memory usage caused by KV caches. \text{BitNet b1.58} marks an important advance toward native long sequence support by reducing activation precision from 16 bits to 8 bits, which effectively doubles the context length under the same resource constraints. This can potentially be further compressed losslessly to 4 bits or even less for 1.58-bit LLMs, which we propose as a direction for future research.

\textbf{LLMs on Edge and Mobile}

Employing 1.58-bit LLMs holds great promise for enhancing language model performance on edge and mobile platforms. These devices are typically constrained by limited memory and computational capacity, which restricts both the model size and performance. The decreased memory footprint and energy usage of 1.58-bit LLMs enable their deployment on such devices, unlocking a variety of applications that were previously unattainable. This advancement can substantially boost the functionality of edge and mobile devices and pave the way for innovative LLM applications. Furthermore, 1.58-bit LLMs are better suited for CPU architectures, which dominate edge and mobile hardware. Consequently, \text{BitNet b1.58} can run efficiently on these platforms, further enhancing their performance and potential.

\textbf{New Hardware for 1-bit LLMs}

Recent efforts such as Groq\footnote{\url{https://groq.com/}} have shown promising outcomes and potential in developing specialized hardware (e.g., LPUs) tailored for LLMs. Building on this progress, we advocate for and foresee the development of novel hardware and systems explicitly optimized for 1-bit LLMs, inspired by the new computational paradigm introduced by BitNet~\cite{bitnet}.